\section{Discussion}\label{sec:discussion}

\subsection{Task Difficulty and the Role of Trivial Baselines}

The trivial baseline analysis (Table~\ref{tab:trivial_baselines}) shows that the three toolpath types are distinguishable from program metadata alone, whether from G-code line range, mean axis position, or even row count. The 100\% RF accuracy achieved by our sensor features, while technically correct, therefore does not by itself demonstrate the value of sensor-based classification or air-cut referencing for this specific three-class task.

We report this finding because it establishes the true difficulty floor for this classification problem and prevents over-claiming. The sensor-based approach remains valuable for three reasons: (1) it works without access to program metadata, which may be unavailable or unreliable in deployment scenarios; (2) it provides the temporal fingerprint representation that offers interpretable physical insight; and (3) it establishes a methodology that transfers to harder discrimination tasks where metadata cannot help. The damaged-spindle detection result (Section~\ref{sec:damage}) demonstrates the third point. There the toolpath strategies are shared, so neither program metadata nor workspace position can separate the classes, yet the sensors reach 98.5\% accuracy using process-dynamics channels rather than the magnetometer position encoder; because the damaged runs are air-cuts, spindle condition and cutting engagement are confounded in that separation, which cannot be attributed to spindle condition alone.

\subsection{Effectiveness of Air-Cut Referencing}

Despite the easy task, three lines of evidence demonstrate a genuine contribution of air-cut referencing.

\textbf{Statistical significance (a consistency check, not the main evidence).} On the full feature set the air-cut effect is small and ceiling-bound. The Friedman test ($\chi^2 = 33.0$, $p < 10^{-6}$, $p_{\mathrm{Bonf}} = 0.011$) is significant, and the bootstrap CIs do not overlap (raw $[0.982, 0.994]$ vs.\ referenced $[1.000, 1.000]$). However, the gap is only $\sim$1~pp at a 100\% ceiling, Cliff's $\delta = 0.022$ is negligible, and the significance traces to a single recurrently-misclassified run. Cohen's $d$ is not meaningful here because the referenced methods have zero variance at the ceiling, so we do not interpret it. We therefore treat the full-feature tests as a consistency check and rest the case for air-cut referencing on the confound-free experiment (raw RF 96.7\% vs.\ subtracted/ratio 100\%) together with the noise-resilience and data-efficiency results below.

\textbf{Robustness.} Across 100 repeated CV evaluations, raw RF features occasionally misclassify (minimum fold accuracy 90\%), while referenced methods maintain perfect accuracy. Per-sample analysis traces this to a single borderline run (pocket150025\_001) that raw RF misclassifies 55\% of the time but all referenced methods classify correctly 100\% of the time.

\textbf{Noise resilience.} Subtracted features maintain 100\% accuracy up to $1.0\times$ additive noise (equal to the feature's natural standard deviation), while raw features degrade to 97.8\% (Figure~\ref{fig:noise_robustness}). At $2.0\times$ noise, the gap widens to 94.6\% vs.\ 86.3\%. This suggests that air-cut referencing would be particularly valuable in production environments with higher sensor noise, calibration drift, or environmental interference.

The two-way ANOVA shows that classifier choice ($\eta^2 = 0.17$) explains approximately $4\times$ more variance than isolation method ($\eta^2 = 0.04$), indicating that for this task, selecting the right classifier matters more than the referencing strategy.

For SVM and MLP classifiers, air-cut referencing provides 4--8 percentage point improvements in mean accuracy and reduces variance, suggesting it removes machine-specific baseline variation that these classifiers are sensitive to.

\subsection{Temporal Structure in Toolpath Fingerprints}

The temporal binning approach outperforms whole-run summary by 2 percentage points (100\% vs.\ 98\%), confirming that temporal structure carries discriminative information. The temporal fingerprint visualizations (Figure~\ref{fig:fingerprint}) reveal characteristic patterns: adaptive clearing produces gradual, continuous signal evolution; facing generates periodic oscillations corresponding to linear passes; and pocketing shows distinct phase transitions from contour passes.

The use of different G-code programs for air and active conditions necessitated temporal binning rather than point-by-point G-code-aligned subtraction. Despite this constraint, the temporal approach captures phase alignment (approach, cutting, retract) through normalized time binning. Future data collection with identical programs would enable a comparison of these approaches.

\subsection{What the Sensors Are Actually Detecting}

The physical interpretation analysis (Section~\ref{sec:physics}) shows that the dominant discriminative features are magnetometer channels that primarily encode \emph{workspace position} (proximity to motors and ferromagnetic structure) rather than \emph{cutting dynamics} (forces, vibrations, torque). This is the sensor-domain equivalent of the trivial position baseline; the classifier can distinguish toolpath types largely because the three programs move through different regions of the machine's workspace.

This observation has three implications:
\begin{itemize}
    \item \textbf{For this study}: the high classification accuracy is partly attributable to program geometry rather than process physics. Future work testing on toolpath variants that traverse the same workspace region would isolate the cutting-dynamics contribution.
    \item \textbf{For sensor selection}: magnetometers on CNC machines are sensitive to position due to motor fields, making them effective for operation identification but poor for process-independent fingerprinting. Accelerometers and current sensors, though less discriminative in this study, capture signals more directly related to cutting mechanics and would likely perform better on within-type discrimination tasks.
    \item \textbf{For air-cut referencing}: the subtracted (active $-$ air) features partially remove position-dependent offsets since the air-cut reference captures the same workspace trajectory. This may explain why the noise robustness analysis (Figure~\ref{fig:noise_robustness}) shows referenced features outperforming raw features, as the subtraction would remove position-correlated noise components.
\end{itemize}

Spindle current, despite being the most direct proxy for cutting torque, shows almost no variation across toolpath types (5.245--5.257~A). This is consistent with the shallow depth of cut (0.025'') producing minimal load relative to the spindle's no-load draw at 12{,}000~RPM. In industrial machines with heavier cuts, spindle current would likely become a more discriminative channel.

\subsection{Sensor Deployment Guidelines}

The ablation study provides guidance for deploying toolpath fingerprinting systems:
\begin{itemize}
    \item A single bed-mounted IMU (y\_bed\_\_3) achieves 100\% accuracy, the best single-sensor placement in this study; we attribute its advantage to proximity to both the workpiece and the Y-axis motor.
    \item Any single IMU module achieves $\geq$94\% accuracy regardless of placement.
    \item Electrical current sensors alone achieve only 76.5\% (the weakest sensor group; an earlier 96.0\% figure was inflated by the spindle2 feature leak, Section~\ref{sec:results}), so IMU-based sensing is preferable where it can be installed.
    \item Two statistics (mean + RMS) per temporal bin achieve 98\%; the full 7-statistic set reaches 100\%.
    \item PCA hurts RF (88.2\% vs.\ 100\%) but helps SVM (94\% at 5 components) and MLP; dimensionality reduction must be matched to the classifier.
\end{itemize}

\subsection{Temporal Drift and Run-Order Effects}

Run-order analysis reveals statistically significant temporal drift in several sensor channels, particularly for the adaptive clearing class. Temperature correlates strongly with run order (Spearman $r = 0.98$, $p < 0.001$), and the bed accelerometer Az shows similar drift ($r = 0.97$). This indicates that the machine warms up over the course of a recording session, producing systematic feature changes correlated with run number rather than toolpath type. Since runs within each toolpath class were collected sequentially, this temperature drift is partially confounded with within-class ordering but does not create between-class confounds (all three types show warming trends).

\subsection{Error Analysis: Borderline Run pocket150025\_001}

Per-sample analysis identified pocket150025\_001 as the single run that raw RF misclassifies in 55\% of repeated CV evaluations. This run deviates from the pocketing class centroid by $> 3.9\sigma$ on 15 features, primarily color sensor variability (ColorR, ColorG, ColorB, ColorA) and frame\_l3 gyroscope statistics. On all 15 deviating features, the sample's value is closer to the facing class centroid than to the pocketing centroid. The color sensor dominance suggests an ambient lighting anomaly during this specific run (e.g., a shadow change). Air-cut referencing resolves this misclassification by normalizing out the session-specific color baseline, demonstrating its value for handling run-level environmental variability.

\subsection{Causal Structure of the Classification Problem}\label{sec:causal_dag}

The causal structure shown in Figure~\ref{fig:causal_dag} explains the empirical findings of this study, namely trivial baselines achieving 100\%, magnetometer dominance via position encoding, and a confound-free pipeline that still benefits from air-cut referencing. The G-code program (an exogenous variable corresponding to the toolpath type) determines both the sequence of motion commands and the workspace region traversed. Workspace region in turn determines magnetometer readings (via proximity to stepper motors and ferromagnetic structure) and motor current dynamics (via different acceleration profiles). Cutting forces are an additional contribution from material removal, which under our shallow-DOC UHMW-PE conditions is small relative to machine-specific signals.

\begin{figure}[H]
\centering
\begin{tikzpicture}[
    node distance=1.4cm and 1.6cm,
    every node/.style={font=\small},
    var/.style={rectangle, draw, rounded corners, minimum width=2.6cm, minimum height=0.8cm, align=center},
    obs/.style={var, fill=blue!10},
    confound/.style={var, fill=orange!15},
    process/.style={var, fill=green!15},
    arrow/.style={-{Latex[length=2mm]}, thick},
    intervene/.style={->, dashed, thick, red},
]
% Top row
\node[var] (gcode) {G-code program\\(toolpath type)};
\node[confound, right=of gcode] (workspace) {Workspace region\\(XYZ trajectory)};
\node[process, right=of workspace] (forces) {Cutting forces\\(material removal)};

% Middle row
\node[obs, below=of workspace] (mag) {Magnetometer\\readings};
\node[obs, below=of forces] (motion) {Accelerometer /\\motor current};

% Bottom row
\node[var, below right=1.0cm and 0.2cm of mag] (features) {Sensor features\\(110 channels)};
\node[var, right=of features] (label) {Predicted\\toolpath class};

% Air-cut reference (off to side)
\node[obs, left=2.0cm of mag] (air) {Air-cut\\baseline\\$\boldsymbol{\mu}_\text{air}$};

% Arrows
\draw[arrow] (gcode) -- (workspace);
\draw[arrow] (gcode) -- (forces);
\draw[arrow] (workspace) -- (mag);
\draw[arrow] (workspace) -- (motion);
\draw[arrow] (forces) -- (motion);
\draw[arrow] (mag) -- (features);
\draw[arrow] (motion) -- (features);
\draw[arrow] (features) -- (label);

% Air-cut intervention
\draw[intervene] (air) to[out=270, in=180] node[below, font=\scriptsize, sloped] {subtract} (features);
\end{tikzpicture}
\caption{Causal structure of the toolpath classification task. Solid arrows represent causal dependencies in the data-generating process. The G-code program is the root cause of both workspace region (a confound) and cutting forces (the process-relevant signal). Magnetometer features are downstream of workspace region only, making them position encoders. Accelerometer and motor-current features are downstream of both workspace region and cutting forces. The dashed red arrow represents the air-cut referencing intervention: subtracting the air-cut baseline from active features removes the workspace-dependent component, leaving the cutting-force component. Because the air-cut and active programs in our dataset traverse \emph{similar but not identical} workspace regions, this subtraction is partial (Section~\ref{sec:results}). Note that under the shallow-DOC UHMW-PE conditions of this study, the cutting-forces $\rightarrow$ motion-features arrow carries a small signal ($\sim$1--2~N at the tool); on metallic workpieces with realistic cutting forces (50--500~N), this arrow would be substantially stronger, shifting the relative balance of position-encoded versus force-encoded discrimination.}
\label{fig:causal_dag}
\end{figure}

This structural model explains four otherwise puzzling observations. First, magnetometer features dominate feature importance because they are directly downstream of the workspace-region confound, which carries nearly all the class signal. Second, removing magnetometer channels does not destroy classification accuracy because accelerometer and motor-current features are also downstream of workspace region (via trajectory-dependent acceleration patterns). Third, air-cut referencing helps even on confound-free features because it removes workspace-dependent components from the accelerometer and current channels as well, isolating their force-relevant residuals. Fourth, cross-condition transfer fails (37\% accuracy) because the air-cut and active-cut conditions used different G-code programs and thus produced different workspace regions, so the position-encoded features do not transfer; the small force-relevant signal would transfer but is dominated by the position component in the raw representation. Together these explain why the paper's headline result (100\% accuracy) is not the contribution. The contribution is identifying the causal structure that makes such results possible and providing a methodology for separating the position-encoded path from the force-relevant path.

\subsection{Limitations and Future Work}

\textbf{No independent test set.} Beyond the single internal held-out split (Section~\ref{sec:results}: 12 of the 51 runs, evaluated once), all evaluations rely on cross-validation (5-fold stratified, 20 repeats) on the same 51 active-cutting runs. That internal split shares the machine, session, and CAM programs of the training data, so it guards against fold-reuse leakage but is \emph{not} an independent test of generalization. With 100 CV evaluations, every sample also appears in many train/test splits. An \emph{independent} test set from a future collection round (a new session, and ideally a second machine) would provide stronger evidence of generalization. Within the existing data, the leave-$K$-runs-out analysis ($K \leq 3$) is the closest approximation.

\textbf{Single machine.} All data is collected from one Bantam Tools Explorer CNC mill. The central hypothesis, that air-cut referencing produces machine-independent fingerprints, cannot be fully validated without cross-machine experiments. The strong generalization results (leave-$K$-out accuracy) suggest robustness to within-machine variation but do not address between-machine variation.

\textbf{Different CAM programs.} The air-cut and active-cutting data used different CAM-generated G-code programs, preventing point-by-point signal subtraction. The temporal binning approach compensates for this, but future data collection should use identical G-code programs for both conditions to enable direct comparison.

\textbf{Limited toolpath diversity.} Three toolpath types with fixed cutting parameters (0.150'' stepover, 0.025'' depth) provide a clean classification task but do not test discrimination between subtle variations within a toolpath type.

\textbf{Workpiece material and cutting forces.} The UHMW-PE workpiece produces extremely low cutting forces ($\sim$1--2~N for a \SI{6.35}{\milli\meter} endmill at \SI{0.635}{\milli\meter} DOC) compared to metals. This means the ``process-specific signal'' isolated by air-cut referencing is dominated by trajectory differences and quasi-static motion dynamics rather than cutting force signatures. On metallic workpieces with higher cutting forces, the sensor signals would contain a larger cutting-force component, potentially making the accelerometer and current channels more discriminative relative to the magnetometer position proxy.

\textbf{Low sampling rate and decimation.} The \SI{4}{\hertz} effective sampling rate (decimated from \SI{119}{\hertz} IMU and \SI{1}{\kilo\hertz} DAQ without explicit anti-aliasing) captures only quasi-static machine motion, not cutting vibration. Higher-rate acquisition retaining the native DAQ bandwidth would enable frequency-domain features (e.g., power spectral density at cutting-relevant frequencies) that are unavailable in this study.

\textbf{Small dataset.} With 51 active-cutting runs and $\leq$19 runs per class, the dataset is adequate for the current three-class problem but may be insufficient for more complex tasks or deep learning approaches.

\textbf{Near-perfect accuracy.} The 100\% RF accuracy across all methods raises concerns about task difficulty. The three toolpath types are structurally quite different (continuous adaptive, back-and-forth facing, concentric pocketing), making them inherently easy to distinguish. Testing on more similar toolpath variants would better evaluate the methods' discriminative power.
